\subsection{Latent Factor and Component Analysis}

Latent factor methods model observed coordinates as mixtures of unobserved sources.  The main question is which structure makes the sources interpretable: covariance structure, statistical independence, or non-Gaussian projections.

\begin{definition} \label{def:sl-factor-analysis-latent-model} \hlt{Factor Analysis}\\
Let $X\in\mathbb{R}^p$ and $q<p$. A \hlt{factor-analysis model} is a representation $X=AS+\varepsilon$ with $A\in\mathbb{R}^{p\times q}$, $S\in\mathbb{R}^q$, $\varepsilon\in\mathbb{R}^p$, where: $\operatorname{Cov}(S)=I_q$, $\operatorname{Cov}(\varepsilon)=D_\varepsilon$ for diagonal $D_\varepsilon\in\mathbb{R}^{p\times p}$, and $\operatorname{Cov}(S,\varepsilon)=0$.\\
The columns $a_1,\ldots,a_q$ of $A$ are the \hlt{factor loadings}, and $\Sigma\defeq\operatorname{Cov}(X)=AA^T+D_\varepsilon$.\\
For any orthogonal $R\in\mathbb{R}^{q\times q}$, $(A,S)$ and $(AR^T,RS)$ yield the same $\Sigma$.\\
Hence $(A,S)$ is identifiable only up to orthogonal rotation.
\end{definition}

\begin{definition} \label{def:sl-independent-component-analysis} \hlt{Independent Component Analysis}\\
Let $X\in\mathbb{R}^p$. An \hlt{independent component model} is a representation $X=AS$ such that:
\begin{enumerate}[label=(\roman*)]
\setlength{\itemsep}{0pt}
\item $A\in\mathbb{R}^{p\times p}$ is nonsingular.
\item $S=(S_1,\ldots,S_p)^T\in\mathbb{R}^p$ has mutually independent components.
\end{enumerate}
If $X$ is whitened, then $\operatorname{Cov}(X)=I_p$ and $A$ may be taken orthogonal.\\
ICA estimates orthogonal $A$ with $Y=A^TX$ approximately componentwise independent, typically via non-Gaussianity contrasts. Gaussian components are identifiable only up to orthogonal rotations.
\end{definition}

\begin{figure}[H]
\centering
\begin{tikzpicture}[
  >=Stealth,
  every node/.style={font=\scriptsize},
  axis/.style={->, gray!65, line width=0.4pt},
  pca/.style={-{Stealth[length=1.6mm]}, red!75!black, line width=0.65pt},
  ica/.style={-{Stealth[length=1.6mm]}, blue!70!black, line width=0.65pt}
]
  \path[use as bounding box] (-1.55,-1.55) rectangle (5.35,1.65);

  \begin{scope}[scale=0.76]
    \draw[axis] (-1.85,0) -- (1.95,0) node[above] {$x_1$};
    \draw[axis] (0,-1.70) -- (0,1.82) node[above] {$x_2$};
    \foreach \x/\y in {
      -1.55/-1.25,-1.30/-0.95,-1.05/-1.42,-0.80/-0.62,-0.55/-1.05,
      -0.35/-0.45,-0.15/-0.85,0.10/-0.28,0.35/-0.72,0.55/-0.10,
      0.80/-0.50,1.05/0.05,1.25/-0.25,1.50/0.35,
      -1.30/0.82,-1.00/1.08,-0.75/0.58,-0.50/1.32,-0.25/0.72,
      0.00/1.05,0.25/0.55,0.50/0.92,0.75/0.40,1.00/0.72,1.30/0.18}
      \fill[blue!65!black] (\x,\y) circle (1.1pt);
    \draw[pca] (0,0) -- (1.48,0.34) node[above right] {$v_1$};
    \draw[pca] (0,0) -- (-0.34,1.48) node[above left] {$v_2$};
  \end{scope}

  \begin{scope}[xshift=3.55cm, scale=0.76]
    \draw[axis] (-1.85,0) -- (1.95,0) node[above] {$y_1$};
    \draw[axis] (0,-1.70) -- (0,1.82) node[above] {$y_2$};
    \foreach \x/\y in {
      -1.70/-0.18,-1.42/0.15,-1.20/-0.05,-0.95/0.18,-0.70/-0.15,
      -0.42/0.08,-0.18/-0.12,0.05/0.10,0.35/-0.08,0.62/0.18,
      0.90/-0.10,1.18/0.08,1.45/-0.15,1.68/0.15,
      -0.15/-1.55,0.18/-1.28,-0.05/-1.02,0.15/-0.72,-0.12/-0.45,
      0.08/-0.15,-0.15/0.12,0.18/0.42,-0.08/0.70,0.12/0.98,
      -0.18/1.25,0.15/1.55}
      \fill[green!45!black] (\x,\y) circle (1.1pt);
    \draw[ica] (0,0) -- (1.55,0) node[below right] {$a_1$};
    \draw[ica] (0,0) -- (0,1.55) node[above left] {$a_2$};
  \end{scope}
\end{tikzpicture}

\caption{\textcolor{red!75!black}{$v_1,v_2$}: PCA variance axes; \textcolor{blue!70!black}{$a_1,a_2$}: ICA rotation exposing independent non-Gaussian coordinates.}
\label{fig:sl-ica-rotation-nongaussianity}
\end{figure}

\begin{method} \label{meth:sl-ica-entropy-view} \hlt{Entropy View of ICA}\\
For a density $g$, define $H(g)=-\int g(y)\log g(y)\,dy$.\\
Let $Y=(Y_1,\ldots,Y_p)$ have joint density $g_Y$ and marginal densities $g_j$.
\begin{enumerate}[label=(\roman*)]
\setlength{\itemsep}{0pt}
\item Set $H(Y)\defeq H(g_Y)$, $H(Y_j)\defeq H(g_j)$, and $I(Y)\defeq\sum_{j=1}^pH(Y_j)-H(Y)$.\\
Then $I(Y)\ge0$, with $I(Y)=0$ iff $Y_1,\ldots,Y_p$ are independent.
\item If $X$ is whitened, $A$ is orthogonal, and $Y=A^TX$, then $H(Y)=H(X)$.\\
Hence minimising $I(Y)$ over orthogonal $A$ is equivalent to minimising $\sum_{j=1}^pH(Y_j)$.
\item Replace marginal-entropy minimisation by marginal non-Gaussianity maximisation,\\
e.g. through negentropy contrasts, since Gaussian marginals maximise entropy at fixed variance.
\end{enumerate}
\end{method}

\begin{remark} \label{rem:sl-projection-pursuit-and-ica} \hlt{Projection Pursuit and ICA}\\
Projection pursuit optimizes a non-Gaussianity index over $Y=a^TX$ or $Y=(a_1^TX,a_2^TX)$.\\
ICA uses the same type of index on $Y_j=a_j^TX$, but imposes an orthogonal system $(a_1,\ldots,a_p)$ after whitening.
\end{remark}

\begin{method} \label{meth:sl-product-density-ica} \hlt{Product-Density ICA}\\
For whitened $X\in\mathbb{R}^p$ and orthogonal $A\in\mathbb{R}^{p\times p}$, set $Y=A^TX$.
\begin{enumerate}[label=(\roman*)]
\setlength{\itemsep}{0pt}
\item Assume the product density $f_Y(y)=\prod_{j=1}^pf_j(y_j)$.
\item Parameterize each marginal density by $f_j(s)=\phi(s)\exp(g_j(s))$ with $\int\phi(s)\exp(g_j(s))\,ds=1$.
\item Estimate $A,g_1,\ldots,g_p$ by penalized likelihood subject to $A^TA=I$.\\
For fixed $A$, the objective decomposes into $p$ one-dimensional smooth density fits.
\end{enumerate}
The roughness penalty shrinks $g_j$ toward $0$, the Gaussian baseline $\phi$.
\end{method}

\begin{breakablealgorithm}
\caption{Product-Density ICA Alternation}
\label{alg:sl-product-density-ica}
\begin{algorithmic}
\Require Whitened observations $x_1,\ldots,x_n$, initial orthogonal matrix $A$
\Ensure Orthogonal unmixing matrix $A$ and components $s_{ij}=a_j^Tx_i$
\Repeat
\For{each component $j$}
\State Fit a smooth tilted-Gaussian density $f_j(s)=\phi(s)\exp(g_j(s))$ to the current scores $a_j^Tx_i$
\State Update $a_j\gets \mathbb{E}\{Xg_j'(a_j^TX)\}-\mathbb{E}\{g_j''(a_j^TX)\}a_j$
\EndFor
\State Orthogonalize $A$: if $A=UDV^T$ is its SVD, set $A\gets UV^T$
\Until{the directions stabilize}
\end{algorithmic}
\end{breakablealgorithm}
